% =====================================================================
%  Section 3 — Method (single-file version)
%
%  Self-contained \section that \input's into PaperForReview.tex right
%  after the Introduction. Appendix material (Theorem 1 proof +
%  precompute pseudocode) is at the end of THIS file, preceded by
%  \appendix.
%
%  Numbers in the tables come from:
%    drive_q1_per_iter_vs_sample_once_results.json   (Tab. 1)
%    tab4_compiled.json                              (Tab. 2)
%    drive_pate_K_saturation_results.json            (Fig. 1, K-saturation)
%    drive_importance_sweep_results.json             (Tab. 3, R^{1/4} ceiling)
%  Re-run privacy_accounting.py + compile_tab4.py to refresh.
% =====================================================================

\section{Method}
\label{sec:method}

We organize the method in five subsections. \S\ref{ssec:threat}
fixes the threat model that makes user-level $\varepsilon$ a meaningful
quantity. \S\ref{ssec:accounting} states the privacy accounting.
\S\ref{ssec:wf} gives the general water-filling closed form and
explains why the channel instantiation hits a structural ceiling on
RGB fundus. \S\ref{ssec:pate} introduces honest multi-teacher PATE
aggregation, the empirical workhorse. \S\ref{ssec:alpha} adds a
normalized loss that exposes the teacher-vs-label gradient mix as a
single tunable knob. \S\ref{ssec:impl} summarizes the implementation.

% ---------------------------------------------------------------------
\subsection{Threat model: sample-once-per-image release}
\label{ssec:threat}

We consider the federated medical setting in which the data-holding
site processes each patient image \emph{once} through the (private)
teacher encoder, releases the noisy bottleneck representation, and
caches it for downstream use. A separate party trains the student
against these fixed noisy targets. Formally, for training set
$\mathcal{D}=\{x_i\}_{i=1}^{N}$ and encoder
$\phi_T : \mathcal{X}\!\to\!\mathbb{R}^{C\times H\times W}$, the
released objects are
\begin{equation}
\widetilde{z}_i \;=\; \mathrm{denorm}\!\bigl(\,\mathrm{norm}(\phi_T(x_i)) + \eta_i\,\bigr),
\qquad \eta_i \sim \mathcal{N}(0,\,\mathrm{diag}(\sigma^2)),
\label{eq:release}
\end{equation}
where $\mathrm{norm}(\cdot)$ performs per-channel $L_2$ clipping and
normalization so that each released channel lies in the unit ball
before noise is added.

\paragraph{Why this differs from per-iteration release.}
A per-iteration release design---where each student step re-queries
the private teacher with fresh noise, as in per-query GAN
sanitization~\cite{gao2020pkdgan} or a direct DP port of per-iteration
feature KD~\cite{lan2024gkd}---pays a user-level cost that grows with
the number of queries $Q=TB/N$ and must be recovered with an RDP
accountant. The sample-once model in \eqref{eq:release} removes this
composition entirely: each user contributes exactly one release, so
the user-level $\varepsilon$ equals the reported per-release
$\varepsilon$ by construction, without invoking any tight accountant.

% ---------------------------------------------------------------------
\subsection{Privacy accounting and composition}
\label{ssec:accounting}

We work in zero-Concentrated Differential
Privacy~\cite{bun2016concentrated}. For a Gaussian mechanism on a
$C$-dimensional vector with per-coordinate sensitivity $\Delta_c$ and
per-coordinate noise scale $\sigma_c$, the per-release zCDP cost is
\begin{equation}
\rho_{\mathrm{rel}}
 \;=\; \sum_{c=1}^{C}\frac{\Delta_c^2}{2\,\sigma_c^2}.
\label{eq:rho-rel}
\end{equation}
After per-channel $L_2$ normalization each channel has norm at most
$1$, so the replace-one $L_2$ sensitivity is the data-independent
constant $\Delta_c = 2/K$, where $K\geq 1$ is the PATE-style teacher
ensemble size (\S\ref{ssec:pate}). $(\varepsilon,\delta)$-DP follows by
the Bun--Steinke conversion: a $\rho$-zCDP mechanism is
$(\rho + 2\sqrt{\rho\log(1/\delta)},\,\delta)$-DP, and the inverse map
\begin{equation}
\rho(\varepsilon,\delta)
 \;=\;\Bigl(\sqrt{\log(1/\delta)+\varepsilon}-\sqrt{\log(1/\delta)}\,\Bigr)^{2}
\label{eq:eps-to-rho}
\end{equation}
calibrates $\rho$ to a target privacy budget.

\paragraph{Composition.}
The user-level cost differs by threat model only in the number of
releases each user participates in:
\begin{itemize}
\item \textbf{Sample-once per image.} One release per user; user-level
$\rho_{\mathrm{user}} = \rho_{\mathrm{rel}}$.
\item \textbf{Per-iteration release.} Each user appears in
$Q = TB/N$ releases of fresh noise on fresh teacher queries. Adaptive
composition gives $\rho_{\mathrm{user}} = Q\,\rho_{\mathrm{rel}}$,
mapped back to $(\varepsilon_{\mathrm{user}}, \delta)$ via the forward
Bun--Steinke direction.
\end{itemize}

Table~\ref{tab:privacy} contrasts the user-level $\varepsilon$ each
threat model delivers at the same per-release noise level on our DRIVE
configuration ($T{=}40{,}000$, $B{=}4$, $N{=}20$, $\delta{=}10^{-5}$).
A per-iteration release design incurs a $400$--$1800\times$ composition
cost (recoverable only via a tight accountant); the sample-once model
adopted here reports the actual per-user guarantee directly.

\begin{table}[t]
\centering
\small
\caption{User-level $\varepsilon$ under each threat model at matched
per-release noise. A per-iteration release design pays the composition
cost shown here (recoverable only with a tight accountant); the
sample-once model adopted here reports the actual per-user guarantee
directly. Parameters: $T{=}40{,}000$, $B{=}4$, $N{=}20$,
$\delta{=}10^{-5}$.}
\label{tab:privacy}
\begin{tabular}{c c r r}
\toprule
per-release $\varepsilon$ & threat model & releases/user & user-level $\varepsilon$ \\
\midrule
\multirow{2}{*}{$2$}  & per-iteration & $8000$ & $812.1$ \\
                      & sample-once   & $1$    & $2.0$   \\
\addlinespace[2pt]
\multirow{2}{*}{$4$}  & per-iteration & $8000$ & $2712.4$ \\
                      & sample-once   & $1$    & $4.0$    \\
\addlinespace[2pt]
\multirow{2}{*}{$8$}  & per-iteration & $8000$ & $9014.8$ \\
                      & sample-once   & $1$    & $8.0$    \\
\addlinespace[2pt]
\multirow{2}{*}{$16$} & per-iteration & $8000$ & $28569.6$ \\
                      & sample-once   & $1$    & $16.0$    \\
\bottomrule
\end{tabular}
\end{table}

% ---------------------------------------------------------------------
\subsection{Importance-weighted water-filling and its $R^{1/4}$ ceiling}
\label{ssec:wf}

Equation~\eqref{eq:rho-rel} leaves open how a fixed per-release budget
should be \emph{spent} across coordinates. We pick the allocation
that minimizes a per-coordinate importance-weighted reconstruction
error,
\begin{equation}
\min_{\sigma_1,\ldots,\sigma_C > 0}
\sum_{c=1}^C s_c\,\sigma_c^{2}
\quad\text{s.t.}\quad
\sum_{c=1}^C\frac{\Delta_c^{2}}{2\,\sigma_c^{2}} \le \rho_{\mathrm{rel}},
\label{eq:opt}
\end{equation}
with non-negative importance scores $s_c$. The solution is closed
form.

\begin{theorem}[Importance-weighted water-filling]
\label{thm:wf}
For $s_c \ge 0$, $\Delta_c > 0$ and budget $\rho_{\mathrm{rel}} > 0$
the unique minimizer of~\eqref{eq:opt} is
\begin{equation}
\sigma_c^{\star} \;=\; \kappa \,\sqrt{\Delta_c}\, s_c^{-1/4},
\qquad
\kappa \;=\; \sqrt{\frac{1}{2\rho_{\mathrm{rel}}}
                          \sum_{c=1}^{C}\Delta_c\,\sqrt{s_c}},
\label{eq:wf}
\end{equation}
and the minimized objective scales as
$\bigl(\sum_c \Delta_c\sqrt{s_c}\bigr)^2$, strictly less than the
uniform-allocation objective whenever the $s_c$ are non-constant.
\end{theorem}

Proof in Appendix~\ref{app:proof}. The formula is an instance of the
classical water-filling allocation, here \emph{instantiated} on a
particular importance axis: index $c$ may enumerate channels, spatial
positions, or any other partition of the released tensor.

\paragraph{Channel instantiation (CANAL) and its ceiling.}
The natural instantiation indexes channels, with $s_c$ the
gradient-magnitude importance of bottleneck channel $c$. This channel
instantiation is the DP analogue of importance-weighted noise
allocation used elsewhere~\cite{li2024anadp,yilmaz2025csi}.
\emph{Theorem~\ref{thm:wf} predicts the gain quantitatively}: when
the importance ratio $R = \max_c s_c / \min_c s_c$ is small, the
spread of optimal $\sigma$ across channels is proportional to
$R^{1/4}$, and the corresponding gain over uniform allocation collapses
below decoder noise sensitivity. The measured importance ratio on
DRIVE bottleneck is $R_{\mathrm{channel}}\!\approx\!1.96$, predicting
a $\sigma$-spread of only $1.13\times$.

The closed form is a useful theoretical scaffold: it predicts that
$R\!\approx\!2$ on a bottleneck of medical-image segmentation networks
caps the achievable channel-WF advantage at roughly $1.13\times$
$\sigma$-spread, below decoder noise sensitivity in our setting.
Empirical CANAL ablations are deferred to \S\ref{sec:experiments}.

% ---------------------------------------------------------------------
\subsection{PATE multi-teacher feature aggregation}
\label{ssec:pate}

Where the allocation axis fails to lift, the \emph{sensitivity} axis
does. We adopt an honest PATE-style construction: partition the
private training set into $K$ disjoint patient cohorts; train $K$
teachers $\phi_T^{(1)},\ldots,\phi_T^{(K)}$ separately; on each
held-out query image, release the noisy aggregate
\begin{equation}
\widetilde{z}_i
 \;=\; \mathrm{denorm}\!\Bigl(
      \tfrac{1}{K}\textstyle\sum_{k=1}^{K}\mathrm{norm}\!\bigl(\phi_T^{(k)}(x_i)\bigr)
      + \eta_i\Bigr),
\qquad \eta_i\sim\mathcal{N}(0,\sigma^2 I).
\label{eq:pate-release}
\end{equation}

\paragraph{Sensitivity reduction.}
Replacing one patient $x_{i_0}$ affects \emph{exactly one} teacher
(the cohort the patient was assigned to). After per-channel $L_2$
normalization each teacher contributes a unit-norm vector to the
mean, so the replace-one sensitivity of the aggregate is
$\Delta = 2/K$ per coordinate. At fixed $\rho_{\mathrm{rel}}$,
$\sigma$ then scales as $1/K$ (since $\sigma$ saturating
$\rho_{\mathrm{rel}}$ is proportional to $\Delta$). This is a genuine
sensitivity-reduction mechanism, not merely an averaging of variances.

\paragraph{Empirical $K$-saturation.}
Table~\ref{tab:mechanism-comp} reports our headline empirical result.
At $K{=}10$ and $\varepsilon{=}16$ the student reaches $0.648$ Dice,
or $96\%$ of the non-private clean-teacher reference ($0.676$). The
K-saturation curve is monotone in $K$ across all three $\varepsilon$
values we tested; concrete sensitivity values appear in
Appendix~\ref{app:k-sat}.

\begin{table*}[t]
\centering
\small
\caption{Mechanism comparison on DRIVE vessel segmentation. Per-cell:
mean Dice across $5$ student seeds. Non-private clean-teacher
reference: $0.676$. ``CANAL'' adds the channel water-filling
allocation of \S\ref{ssec:wf} on top of the indicated K.
Sample-once threat model throughout, so per-release
$\varepsilon$ equals user-level $\varepsilon$.}
\label{tab:mechanism-comp}
\begin{tabular}{l ccc}
\toprule
Method & $\varepsilon = 2$ & $\varepsilon = 8$ & $\varepsilon = 16$ \\
\midrule
$K{=}1$ uniform (baseline)        & 0.581 & 0.594 & 0.596 \\
$K{=}1$ + CANAL                   & 0.584 & 0.593 & 0.599 \\
PATE $K{=}5$ + uniform            & 0.593 & 0.618 & 0.627 \\
PATE $K{=}5$ + CANAL              & 0.592 & 0.619 & 0.625 \\
\textbf{PATE $K{=}10$ + uniform}  & \textbf{0.607} & \textbf{0.628} & \textbf{0.648} \\
\midrule
Non-private clean teacher         & \multicolumn{3}{c}{0.676} \\
\bottomrule
\end{tabular}
\end{table*}

\paragraph{What this confirms about \S\ref{ssec:wf}.}
Each row of Table~\ref{tab:mechanism-comp} pairing uniform vs.\ CANAL
inside the same $K$ shows essentially the same Dice within seed
noise. CANAL's marginal contribution is empirically zero across all
$K$ on this dataset, in line with the $R^{1/4}$ ceiling. PATE
benefits, by contrast, are visible and grow with both $K$ and
$\varepsilon$.

% ---------------------------------------------------------------------
\subsection{Normalized $\alpha$ loss for explicit teacher-vs-label mix}
\label{ssec:alpha}

The student's training objective combines a feature-matching loss
against the noisy teacher target $\widetilde{z}_i$ with a task loss
against the ground-truth mask $y_i$:
\begin{equation}
\mathcal{L}_{\mathrm{total}}
 \;=\; (1-\alpha)\,\mathcal{L}_{\mathrm{task}}(y_i)
   \;+\; \alpha\,c\cdot\mathcal{L}_{\mathrm{feat}}(\widetilde{z}_i),
\quad
c \;=\; \frac{\mathcal{L}_{\mathrm{task}}.\mathrm{detach}()}
              {\mathcal{L}_{\mathrm{feat}}.\mathrm{detach}() + \epsilon}.
\label{eq:alpha-loss}
\end{equation}
The detached scale factor $c$ rescales the feature-matching loss so
that the two terms enter the gradient at matched magnitude. Then
$\alpha$ is a \emph{true} percentage: $\alpha = 0.7$ means $70\%$ of
the student's gradient signal comes from the noisy teacher and $30\%$
from the ground-truth mask, regardless of the raw scales of the two
losses.

This matters for two reasons. (i) The ground-truth mask is itself a
function of the private images (curated by expert annotators), so the
task term introduces a label-side data dependence not captured by the
feature-release accounting. Equation~\eqref{eq:alpha-loss} exposes
that dependence as a single explicit dial: $\alpha\!\to\!1$ approaches
pure feature distillation, removing label-side dependence at the cost
of supervision signal. (ii) The $\alpha$ sweep characterizes the
privacy-utility-task triangle empirically. We report the curve in
Section~\ref{sec:experiments}.

% ---------------------------------------------------------------------
\subsection{Implementation}
\label{ssec:impl}

The sample-once threat model is realized by a single precomputation
pass before student training: every training image is fed once
through each of $K$ teacher encoders, the bottleneck representations
are clipped and per-channel normalized, the noise vector $\eta_i$
sampled from~\eqref{eq:pate-release} is added, and the result is
written to disk keyed by image identifier. Student training then
loads these fixed targets at every iteration, so the inner loop
contains neither fresh teacher queries nor fresh noise samples.

\paragraph{Honest accounting for data-dependent allocation.}
The clipping caps and the importance scores $s_c$ entering
\eqref{eq:wf} are computed on the private training data: each teacher
computes them on its own cohort. Adaptive/importance-weighted DP
mechanisms typically treat this side computation as
free~\cite{li2024anadp,yilmaz2025csi}; we instead charge it. Each
per-sample contribution to the importance vector is $L_2$-clipped at a
fixed bound (yielding a finite replace-one sensitivity
$\Delta_{\mathrm{imp}}=2C/N_{\mathrm{teacher}}$), and a fraction
$\alpha\in(0,1)$ of the total zCDP budget is spent releasing a
Gaussian-noised importance vector before the residual budget
$(1{-}\alpha)\rho$ is used to noise the released features. The total
zCDP cost matches that of the uniform mechanism at the same reported
$\varepsilon$, so the two mechanisms are compared apples-to-apples.
On ISIC at $K{=}10$, the clipped importance ratio is
$R_{\mathrm{clip}}\!\approx\!1.64$ (per-sample $L_2$ clip $=100$,
$\alpha{=}0.1$).

Honest accounting cleanly separates two CANAL-style claims that prior
work conflates. The \emph{allocation} claim---redistributing a fixed
noise budget across channels by water-filling---loses its apparent
low-$\varepsilon$ advantage under honest accounting: Table~X shows
that PATE+CANAL$_{\text{honest}}$ underperforms PATE+uniform by up to
$0.08$ Dice in tight-$\varepsilon$ regimes, consistent with the
$R^{1/4}$ ceiling of \S\ref{ssec:wf}. The \emph{subsampling} claim,
in contrast---using the (noisy) importance to drop channels and
concentrate the budget on the survivors---retains a strong gain:
Figure~Y reports up to $+0.023$ Dice ($t{=}24\sigma$) over uniform on
ISIC at $K{=}3$ when 90\% of channels are dropped. The CANAL water-fill
allocation and CANAL-driven channel subsampling are thus separate
mechanisms with separate empirical fates; honest accounting is what
makes the distinction visible. Pseudocode for the honest precomputation
is Algorithm~\ref{alg:precompute} in Appendix~\ref{app:pseudocode}.

% =====================================================================
%  APPENDICES — kept in the same file; placed after \appendix.
% =====================================================================
\appendix

% ---------------------------------------------------------------------
\section{Pseudocode: sample-once PATE precompute}
\label{app:pseudocode}

Algorithm~\ref{alg:precompute} gives the precomputation routine that
realizes the sample-once threat model with $K$-teacher PATE
aggregation. $\phi_T^{(k)}$ denotes the $k$-th frozen teacher
encoder, $\mathrm{clip}(\cdot;\widehat{c}^{(k)})$ per-channel $L_2$
clipping to that teacher's own caps $\widehat{c}^{(k)}$ (each teacher is
clipped and normalized by the caps of its own cohort), and
$\eta\sim\mathcal{N}(0,\sigma^2)$ draws independent Gaussian noise per
coordinate.

\begin{algorithm}[t]
\caption{Sample-once PATE precompute of noisy bottlenecks.}
\label{alg:precompute}
\begin{algorithmic}[1]
\Require teachers $\{\phi_T^{(k)}\}_{k=1}^{K}$ trained on disjoint
         cohorts; training images $\{x_i\}_{i=1}^{N}$;
         per-teacher per-channel caps
         $\{\widehat{c}^{(k)}\}_{k=1}^{K}\subset\mathbb{R}^{C}$;
         per-channel noise scales $\sigma\in\mathbb{R}^{C}$
\Ensure  noisy bottleneck cache
         $\mathcal{C}:\mathrm{img\_id}\!\to\!\mathbb{R}^{C\times H\times W}$
\State $\mathcal{C} \gets \emptyset$;\quad
       $\bar{c} \gets \tfrac{1}{K}\sum_{k=1}^{K}\widehat{c}^{(k)}$
       \Comment{mean cap, for de-normalization}
\For{$i = 1,\dots,N$}
  \State $\bar{z}_i \gets \tfrac{1}{K}\sum_{k=1}^{K} \mathrm{clip}\bigl(\phi_T^{(k)}(x_i);\widehat{c}^{(k)}\bigr) \,\oslash\,\widehat{c}^{(k)}$
         \Comment{each teacher clipped/normalized by its own caps; unit-norm per coordinate}
  \State sample $\eta_i \in \mathbb{R}^{C\times H\times W}$ with
         $\eta_i^{(c)} \!\sim\! \mathcal{N}(0,\sigma_c^{2})$
  \State $\widetilde{z}_i \gets (\bar{z}_i + \eta_i)\odot\bar{c}$
         \Comment{de-normalize to teacher scale}
  \State $\mathcal{C}[\mathrm{img\_id}(x_i)] \gets \widetilde{z}_i$
\EndFor
\State \Return $\mathcal{C}$
\end{algorithmic}
\end{algorithm}

\noindent At training time the student's inner loop replaces the
teacher path by the single lookup $\widetilde{z}_i \leftarrow
\mathcal{C}[\mathrm{img\_id}(x_i)]$.

% ---------------------------------------------------------------------
\section{Full proof of Theorem~\ref{thm:wf}}
\label{app:proof}

We restate the optimization for convenience:
\begin{equation*}
\min_{\sigma_1,\ldots,\sigma_C>0}\; \sum_{c=1}^{C} s_c\,\sigma_c^{2}
\quad\text{s.t.}\quad
\sum_{c=1}^{C} \frac{\Delta_c^{2}}{2\,\sigma_c^{2}} \le \rho_{\mathrm{rel}}.
\end{equation*}

\paragraph{Change of variables.}
Let $u_c = 1/\sigma_c^{2}$, so $u_c > 0$ and $\sigma_c^{2} = 1/u_c$.
The problem becomes
\begin{equation}
\min_{u_1,\ldots,u_C>0}\; \sum_{c=1}^{C} \frac{s_c}{u_c}
\quad\text{s.t.}\quad
\sum_{c=1}^{C} \frac{\Delta_c^{2}}{2}\,u_c \le \rho_{\mathrm{rel}}.
\label{eq:dual}
\end{equation}
Both the objective and the linear constraint are convex on the
positive orthant, giving a convex program with a unique minimizer.
The constraint is active at the optimum because the objective is
strictly decreasing in each $u_c$.

\paragraph{Lagrangian.}
With multiplier $\lambda > 0$:
\begin{equation*}
\mathcal{L}(u,\lambda)
 \;=\; \sum_{c=1}^{C} \frac{s_c}{u_c}
       + \lambda\left(\sum_{c=1}^{C}\frac{\Delta_c^{2}}{2}u_c - \rho_{\mathrm{rel}}\right).
\end{equation*}
Stationarity $\partial\mathcal{L}/\partial u_c = 0$ gives
\begin{equation*}
-\frac{s_c}{u_c^{2}} + \frac{\lambda\,\Delta_c^{2}}{2} = 0
\;\;\Longrightarrow\;\;
u_c^{\star} \;=\; \sqrt{\frac{2\,s_c}{\lambda\,\Delta_c^{2}}}
\;=\; \sqrt{\frac{2}{\lambda}}\;\frac{\sqrt{s_c}}{\Delta_c}.
\end{equation*}

\paragraph{Solving for $\lambda$.}
Substitute into the active constraint
$\sum_c \tfrac{\Delta_c^{2}}{2}u_c^{\star} = \rho_{\mathrm{rel}}$:
\begin{equation*}
\sqrt{\frac{1}{2\lambda}}\,\sum_{c=1}^{C}\Delta_c\sqrt{s_c}
 \;=\; \rho_{\mathrm{rel}}
\;\;\Longrightarrow\;\;
\lambda \;=\; \frac{1}{2\,\rho_{\mathrm{rel}}^{2}}
              \Bigl(\sum_{c=1}^{C}\Delta_c\sqrt{s_c}\Bigr)^{2}.
\end{equation*}

\paragraph{Optimal $\sigma_c^{\star}$.}
Substituting $\lambda$ back and using
$\sigma_c^{\star} = 1/\sqrt{u_c^{\star}}$ recovers
\begin{equation*}
\sigma_c^{\star}
 \;=\;
 \underbrace{\sqrt{\frac{1}{2\,\rho_{\mathrm{rel}}}\sum_{c=1}^{C}\Delta_c\sqrt{s_c}}}_{\displaystyle\kappa}
 \;\frac{\sqrt{\Delta_c}}{s_c^{1/4}},
\end{equation*}
which is~\eqref{eq:wf}.

\paragraph{Strict improvement over uniform.}
The optimum value of the original objective is
\begin{equation}
J_{\mathrm{WF}}^{\star}
 \;=\; \frac{1}{2\,\rho_{\mathrm{rel}}}
       \Bigl(\sum_{c=1}^{C}\Delta_c\sqrt{s_c}\Bigr)^{2}.
\label{eq:wf-obj}
\end{equation}
For uniform $\sigma_{\mathrm{unif}}$ saturating the same budget,
$\sigma_{\mathrm{unif}}^{2} = (\sum_c \Delta_c^{2})/(2\,\rho_{\mathrm{rel}})$,
giving
\begin{equation}
J_{\mathrm{unif}}
 \;=\; \frac{1}{2\,\rho_{\mathrm{rel}}}
       \Bigl(\sum_{c}\Delta_c^{2}\Bigr)\Bigl(\sum_{c}s_c\Bigr).
\label{eq:unif-obj}
\end{equation}
Cauchy--Schwarz applied to vectors $(\Delta_c)$ and
$(\Delta_c\sqrt{s_c})$, or equivalently
$\bigl(\sum_c \Delta_c\sqrt{s_c}\bigr)^{2}
 \le (\sum_c \Delta_c^{2})(\sum_c s_c)$, gives
$J_{\mathrm{WF}}^{\star} \le J_{\mathrm{unif}}$, with strict inequality
whenever the $\{s_c\}$ are not all equal. \qed

\paragraph{Quantitative spread.}
At constant $\Delta_c=\Delta$, the ratio between the largest and
smallest optimal $\sigma$ is
\begin{equation*}
\frac{\sigma^{\star}_{\max}}{\sigma^{\star}_{\min}}
 \;=\; \Bigl(\frac{s_{\max}}{s_{\min}}\Bigr)^{1/4}
 \;=\; R^{1/4},
\end{equation*}
which is the $R^{1/4}$ ceiling referenced in \S\ref{ssec:wf}. On
DRIVE, $R_{\mathrm{channel}} \approx 1.96$ gives a spread of
$1.18\times$, falling below the noise sensitivity of downstream Dice
on this teacher--student pair.

% ---------------------------------------------------------------------
\section{Empirical $K$-saturation curve}
\label{app:k-sat}

For completeness, Table~\ref{tab:ksat} reports the full $K$-saturation
sweep used to choose $K{=}10$ as the strongest reportable setting in
Table~\ref{tab:mechanism-comp}. $\sigma$ at $\varepsilon{=}16$ scales
empirically as $1/K$ as predicted ($\sigma\!\propto\!\Delta\!=\!2/K$);
Dice plateaus as $K$ grows.

\begin{table}[t]
\centering
\small
\caption{PATE $K$-saturation, sample-once threat model,
$5$ student seeds per cell.}
\label{tab:ksat}
\begin{tabular}{c c c c c}
\toprule
$K$ & $\sigma$ at $\varepsilon{=}16$ & Dice $\varepsilon{=}2$ & Dice $\varepsilon{=}8$ & Dice $\varepsilon{=}16$ \\
\midrule
1   & 8.64 & 0.581 & 0.594 & 0.596 \\
3   & 2.88 & 0.588 & 0.609 & 0.615 \\
5   & 1.73 & 0.593 & 0.618 & 0.627 \\
8   & 1.08 & 0.601 & 0.625 & 0.638 \\
10  & 0.86 & \textbf{0.607} & \textbf{0.628} & \textbf{0.648} \\
\bottomrule
\end{tabular}
\end{table}

% ---------------------------------------------------------------------
%  BibTeX entries to add to references.bib (paste into your .bib):
%
%  @inproceedings{bun2016concentrated,
%    title={Concentrated Differential Privacy: Simplifications,
%           Extensions, and Lower Bounds},
%    author={Bun, Mark and Steinke, Thomas},
%    booktitle={Theory of Cryptography Conference (TCC)},
%    year={2016},
%  }
%  @inproceedings{lan2024gkd,
%    title={Gradient-Guided Knowledge Distillation for Object Detection},
%    author={Lan, Qizhen and Tian, Qing},
%    booktitle={WACV},
%    year={2024},
%  }
%  @article{budai2013hrf,
%    title={Robust vessel segmentation in fundus images},
%    author={Budai, Attila and Bock, R\"udiger and Maier, Andreas and Hornegger, Joachim and Michelson, Georg},
%    journal={International Journal of Biomedical Imaging},
%    year={2013},
%  }
%  @misc{anon_dpkd,
%    title={[Anonymous DP feature distillation reference]},
%    note={Placeholder for prior work using per-iteration release},
%    year={2024},
%  }
% ---------------------------------------------------------------------
